\chapter{Requirement Analysis}
\label{chap:requirements}

This chapter specifies what BiasAperture must do, its \glspl{fr}, and how well it must do it, its \glspl{nfr}. Hardware and software requirements are given as \emph{system requirements}: what platform and environment are needed to build and run the system. \Cref{chap:methodology} explains \emph{how} these requirements are implemented and validated.

\section{Functional Requirements}

Functional requirements are listed using IEEE~830-style ``shall'' statements, one capability per requirement.

\begin{enumerate}
    \item \textbf{FR-001 (Data Ingestion):} The system shall load, validate, and standardise a benchmark dataset (FairFace or UTKFace) or a compatible custom dataset, checking image integrity and aligning demographic annotation fields (race or ethnicity, perceived gender, age group) into a common internal schema.
    \item \textbf{FR-002 (Dual-Mode Model Interface):} The system shall obtain predictions from a target facial-analysis model either by direct in-process inference against a supplied PyTorch or TensorFlow model object, or by batch ingestion of a precomputed predictions file in \gls{csv} or \gls{json} format.
    \item \textbf{FR-003 (Fairness Metric Computation):} The system shall compute, per subgroup and per demographic intersection, the four Core~Four disparity metrics, \gls{dpd}, \gls{eod}, \gls{eop}, and \gls{dir}, using AIF360 and Fairlearn as independent, cross-validating computational backends.
    \item \textbf{FR-004 (Statistical Significance Testing):} The system shall accompany every reported disparity with a chi-squared test of independence between the predicted-label distribution and the demographic subgroup, and with a bootstrap confidence interval on the subgroup accuracy difference.
    \item \textbf{FR-005 (Explainability):} The system shall generate \gls{shap}-based feature-attribution output identifying which input features are most associated with a detected disparity, to support diagnosis of proxy-variable entanglement.
    \item \textbf{FR-006 (Report Generation):} The system shall assemble the computed metrics, statistical tests, and dataset and model provenance information into a standardised, exportable \gls{html} fairness report, structurally informed by the Model Cards and Datasheets for Datasets conventions, with every metric row showing subgroup sample size, point estimate, 95\% confidence interval, and either a $p$-value or an insufficient-sample flag.
    \item \textbf{FR-007 (Regulatory Traceability):} The system shall map every computed metric, in the generated report, to its corresponding clause under Article~10 or Annex~IV of the \gls{eu} \gls{ai} Act and to its corresponding function within the \gls{nist} \gls{ai} \gls{rmf}.
    \item \textbf{FR-008 (Orchestration and Configuration):} The system shall expose a command-line configuration interface through which an auditor selects the benchmark dataset, the model or predictions source, and the metrics to compute, and shall re-run identically given the same configuration and inputs.
    \item \textbf{FR-009 (Licensing Acknowledgement):} The system shall display and require explicit acknowledgement of the selected dataset's licensing terms before an audit is executed.
\end{enumerate}

\section{Non-Functional Requirements}

\begin{enumerate}
    \item \textbf{NFR-001 (Statistical Rigour):} All significance testing shall use a threshold of \gls{symb:alpha} $= 0.05$, and the system shall report exact $p$-values rather than only a pass/fail flag.
    \item \textbf{NFR-002 (Uncertainty Quantification):} Every reported subgroup accuracy difference shall include a 95\% confidence interval computed from a minimum of 1{,}000 bootstrap resamples.
    \item \textbf{NFR-003 (Data-Integrity Guard):} Any subgroup with a sample size below $n = 30$ shall be automatically flagged as ``insufficient sample, not reported'' rather than assigned a computed metric value.
    \item \textbf{NFR-004 (Performance):} A full-dataset evaluation over FairFace's 108{,}501 images shall complete within 4 hours on a single mid-range \gls{gpu}; a stratified development subset of $n = 5{,}000$ images shall complete within 30 minutes on \gls{cpu} alone.
    \item \textbf{NFR-005 (Modularity):} The system shall maintain strict separation of concerns across its architectural modules, each exposing a narrow, well-defined interface, so that any module can be unit-tested and extended independently of the others.
    \item \textbf{NFR-006 (Reproducibility):} Bootstrap resampling shall use fixed random seeds, and third-party dependency versions (AIF360, Fairlearn, PyTorch or TensorFlow) shall be pinned in a lock-file to prevent silent behavioural drift between runs.
    \item \textbf{NFR-007 (Portability):} The system shall run cross-platform on any environment supporting the pinned Python dependency stack, without requiring proprietary or platform-specific components.
    \item \textbf{NFR-008 (Explainability Performance):} \gls{shap} attribution shall be computed only on flagged disparities rather than the full dataset by default, to keep explainability overhead a small fraction of total runtime.
\end{enumerate}

\section{System Requirements}
This section answers what machine and software stack are required to build and run BiasAperture.

\subsection{Hardware Requirements}
\begin{itemize}
    \item \textbf{Minimum (development, stratified subset):} 64-bit \gls{cpu}, 4 or more cores, 8~GB \gls{ram}, 5~GB free disk space.
    \item \textbf{Recommended (full-dataset audit, NFR-004 target):} A single mid-range \gls{gpu} with 8~GB or more \gls{vram}, for example an NVIDIA T4-class device, either a local workstation \gls{gpu} or a free-tier cloud notebook \gls{gpu} such as Google Colab.
    \item No specialised or custom hardware is required; the platform performs inference and metric computation only, never model training from scratch.
\end{itemize}

\subsection{Software Requirements}
\begin{itemize}
    \item \textbf{Operating system:} Linux, Windows, or macOS; the stack is pure Python and cross-platform.
    \item \textbf{Language and runtime:} Python~3.10 or later.
    \item \textbf{Model interface:} PyTorch or TensorFlow, whichever framework the audited model was built in.
    \item \textbf{Image preprocessing:} OpenCV, pandas, NumPy.
    \item \textbf{Fairness computation:} AIF360, Fairlearn, scikit-learn, SciPy.
    \item \textbf{Explainability:} \gls{shap}.
    \item \textbf{Report generation:} Jinja2 for \gls{html} templating.
\end{itemize}
